\subsection{Spazi di Hilbert e trasformazioni}
Uno spazio di Hilber è uno \textbf{Spazio Vettoriale Complesso} dotato di un prodotto interno (su cui si andrà a definire una norma), le cui successioni di vettori di cauchy convergono asintoticamente a un elemento dello spazio stesso. Tale convergenza, per spazi vettoriali a dimensione finita (come nei nostri casi), è sempre verificata. Tale spazio è spannato su una base ortonormale complessa finita. Le trasformazioni quantistiche saranno viste come matrici unitarie (\(U^\dagger U = U U^\dagger = I, U^\dagger = (U^T)^*\)) moltiplicate per i vettori di stato.

\begin{info}
Dato uno spazio vettoriale complesso \(\mathbb{V}\) dotato di prodotto scalare \(<-,->\), che ne definisce la norma: \(| \cdot | = \sqrt{<\cdot,\cdot>}\), e la distanza: \(d(V_1,V_2) = |V_1 - V_2|\).

Sì definisce successione di Cauchy, una successione di vettori: \(V_0, V_1, V_2, \dots\) tale per cui, \(\forall \epsilon \in \mathbb{R}, \epsilon > 0\)

\[
\exists N_0 \in \mathbb{N} : \forall m,n > N_0, d(V_m, V_n) \le \epsilon
\]

Uno spazio vettoriale prehilbertiano (dotato di prodotto interno) è detto \textbf{completo} se per ogni successione di cauchy (\(V_0, V_1, \dots\)), esiste un vettore \(V \in \mathbb{V}\) tale per cui:
\[
\lim_{n \to \infty} |V_n - V| = 0
\]
Quindi se tali successioni convergono in un elemento interno allo spazio stesso. Sì definisce, quindi, uno \textbf{Spazio di Hilbert}, uno spazio normato e completo.

Ogni spazio vettoriale complesso a dimensione finita, è automaticamente uno spazio completo. Quindi, ogni spazio vettoriale complesso normato a dimensione finita è uno spazio di Hilbert
\end{info}